\section{Balance-of-payments booking of the adoption flows}
\label{app:booking}

The sign and magnitude of the adoption channel's contribution to output depend
on how the green-margin resource flows, the switching cost $\psi_g D^{sw}$ and
the green energy bill, are booked in the balance of payments. The main text
reports the pure-import baseline throughout. This appendix collects the two
results that depend on that choice: the import-content sweep of the adoption
spending, and the carbon-revenue recycling experiment, which is a
booking-conditional statement about where the switching subsidy goes.

\subsection{Import content of the adoption spending}\label{app:booking_flows}
The adoption cost enters the household budget as a CES bundle of home and
foreign goods with import share $\alpha_F^{switch}$. At $\alpha_F^{switch}=1$ the
switching cost is spent entirely on imports, the pure-import baseline of the main
text; lowering it routes adoption spending onto the home good. The share is the
continuous form of the import-versus-domestic booking of
Section~\ref{sec:bop}: the domestic booking of earlier drafts is the
$\alpha_F^{switch}=0$ endpoint, where neither the switching cost nor the green
energy bill is a leakage.

The import content governs the sign of the adoption channel's cumulative output
contribution, $\Delta_{24}$ (full economy minus the frozen-choice counterfactual
at the common steady state). Under the price shock it rises monotonically from
$-9.9$ at $\alpha_F^{switch}=1$, through $+6.3$ at $\alpha_F^{switch}=0.5$, to
$+23.0$ at $\alpha_F^{switch}=0$: the leakage shrinks as spending shifts to the
home good, and the channel turns expansionary once enough of it falls on the home
good. Table~\ref{tab:afs_sweep} reports the grid for both shocks and
Figure~\ref{fig:afs_sweep} the output paths.

Two invariances make clear that this is a booking statement rather than a statement
about the crisis. Peak output stays negative across the whole grid, from
$-2.87\%$ at $\alpha_F^{switch}=1$ to $-1.54\%$ at $\alpha_F^{switch}=0$: the
import content dampens the recession; it never turns it into a boom. The peak
green-share response is flat at roughly $8.0$--$8.2$ pp throughout, so the
adoption margin itself is untouched by where its financing is booked. The policy
ordering of Sections~\ref{sec:shield}--\ref{sec:carbontax} is likewise
unaffected. These invariances are why the paper states its main result in terms of the
margin's response to policy rather than the output sign of the channel.

\begin{table}[H]\centering
\figtitle{Adoption-spending import content}
\IfFileExists{output/tab_afs_sweep.tex}{\input{output/tab_afs_sweep.tex}}{%
\emph{Run \texttt{run\_afs\_sweep.py} to generate \texttt{tab\_afs\_sweep.tex}.}}
\label{tab:afs_sweep}
\fignotes{Both shocks, across the import share $\alpha_F^{switch}$ of the
switching-cost bundle. Columns: peak output, peak green-share response, peak
net-exports-to-GDP, and the adoption channel's cumulative output contribution
$\Delta_{24}$ (adoption$\to y$).}
\end{table}

\begin{figure}[H]\centering
\figtitle{Output response across the adoption-spending import share}
\includegraphics[width=\textwidth]{output/fig_afs_sweep.png}
\fignotes{Import share $\alpha_F^{switch}$ of the switching-cost bundle, price
shock (left) and supply shock (right).}
\label{fig:afs_sweep}
\end{figure}

\subsection{Recycling the carbon revenue: rebate versus green subsidy}
\label{sec:recycle}

The permanent carbon tax of Section~\ref{sec:carbontax} returns its
revenue lump sum. An alternative spends it on the transition itself, paying a
share $s_g$ of the green switching cost. We compare the two recyclings across
carbon-tax rates $\tau_b\in\{0.10,0.20,0.30\}$, holding $\psi_g$ at its no-tax
value so the steady-state green share floats (Figures~\ref{fig:taub_rebate}
and~\ref{fig:taub_green}). Under the lump-sum rebate the steady-state green
share rises from $8.2\%$ to $16.8\%$ across the grid, and the whole brown-price
path in the crisis shifts up by the factor $(1+\tau_b)$: the pre-tax price
passes the world shock through identically, while households pay the
tax-inclusive price. Under \emph{budget-neutral} green-subsidy recycling, $s_g$
solved so that the switching subsidy exactly exhausts the carbon revenue
($T^{\mathrm{reb}}=0$), the revenue funds a subsidy covering $11$--$18\%$ of
the switching cost and lifts the steady-state green share to $68$--$78\%$.

Two features stand out. First, the carbon base is self-eroding: at
$\tau_b=0.10$ the revenue is a third smaller under green-subsidy recycling than
under the rebate, because the greener economy holds less brown-durable stock to
tax. Second, the crisis-time greening response is \emph{smaller} from the greener
starting point under the subsidy, not larger: recycling that greens the
economy through a lower switching \emph{cost} preferentially converts the
households nearest the logit margin, leaving a residual brown pool that is less
responsive to the shock. The rebate does the reverse: it greens
through a higher permanent brown \emph{price} and leaves the residual pool closer
to the margin. It is the adoption-absorber exhaustion of
Section~\ref{sec:routeA} seen across policy instruments rather than across shock
persistence.

\begin{figure}[H]\centering
\figtitle{Carbon-tax sweep: lump-sum rebate}
\includegraphics[width=\textwidth]{output/fig_exante_taub_rebate.png}
\label{fig:taub_rebate}
\fignotes{Laissez-faire and ex-post cap against the ex-ante ETS economy at
$\tau_b\in\{0.10,0.20,0.30\}$, carbon revenue rebated lump sum.}
\end{figure}

\begin{figure}[H]\centering
\figtitle{Carbon-tax sweep: budget-neutral green-subsidy recycling}
\includegraphics[width=\textwidth]{output/fig_exante_taub_greensub_bn.png}
\label{fig:taub_green}
\fignotes{Budget-neutral recycling ($T^{\mathrm{reb}}=0$, $s_g$ endogenous). Same
axes as Figure~\ref{fig:taub_rebate}.}
\end{figure}
